% Compile with: uplatex file.tex && dvipdfmx file.dvi
%           or: latexmk -uplatex file.tex
%\documentclass[uplatex,dvipdfmx]{jsarticle}
\documentclass[dvipdfmx,a4paper,11pt]{jsarticle}
\usepackage{amsmath, amssymb}
\usepackage{tikz}

\title{scenario\_utility\_maximization}
\date{\today}

\begin{document}
\maketitle

コブ・ダグラス型効用関数 $U(x_1,x_2)=x_1^ax_2^b$ を予算制約 $p_1x_1+p_2x_2=m$ の下で最大化する。まず添字つきの記号を「名前」として宣言する——宣言しないと操作時に成分インデックス $x(1)$ として解決され、一階条件が誤る。

$x_1$ をスカラーとする

\[
  x_1
\]

$x_2$ をスカラーとする

\[
  x_2
\]

$p_1$ をスカラーとする

\[
  p_1
\]

$p_2$ をスカラーとする

\[
  p_2
\]

効用関数を登録する。

$U = x_1^a x_2^b$

\begin{equation}
  U = x_1^a x_2^b
  \label{eq:E1}
\end{equation}

ラグランジアンを研究者が書く。乗数の名前も研究者が付ける（システムが自動採番する必要はない）。

$L = U - \lambda (p_1 x_1 + p_2 x_2 - m)$

\begin{equation}
  L = U - \lambda (p_1 x_1 + p_2 x_2 - m)
  \label{eq:E2}
\end{equation}

一階条件のために、ラグランジアンを各変数で偏微分する。

E2 を $x_1$ で偏微分して

\begin{equation}
  \frac{\partial L}{\partial x_1} = x_{1}^{a-1} x_{2}^{b} a-p_{1} \lambda
  \label{eq:E3}
\end{equation}

E2 を $x_2$ で偏微分して

\begin{equation}
  \frac{\partial L}{\partial x_2} = x_{1}^{a} x_{2}^{b-1} b-p_{2} \lambda
  \label{eq:E4}
\end{equation}

E2 を $\lambda$ で偏微分して

\begin{equation}
  \frac{\partial L}{\partial \lambda} = m-p_{2} x_{2}-p_{1} x_{1}
  \label{eq:E5}
\end{equation}

それぞれの右辺を 0 とおいて、一階条件（連立方程式）にする。

E3 の右辺を 0 とする

\begin{equation}
  x_{1}^{a-1} x_{2}^{b} a-p_{1} \lambda = 0
  \label{eq:E6}
\end{equation}

E4 の右辺を 0 とする

\begin{equation}
  x_{1}^{a} x_{2}^{b-1} b-p_{2} \lambda = 0
  \label{eq:E7}
\end{equation}

E5 の右辺を 0 とする

\begin{equation}
  m-p_{2} x_{2}-p_{1} x_{1} = 0
  \label{eq:E8}
\end{equation}

比較静学を先に行う。需要関数を陽に $F(x_1,m)=am-(a+b)p_1x_1=0$ の形で書く。この段は連立を解く前に置く必要がある——解くと $x_1$ が CAS 側で束縛され、$F$ が恒等的に 0 に潰れるためである。

$F = a m - (a+b) p_1 x_1$

\begin{equation}
  F = a m - (a+b) p_1 x_1
  \label{eq:E9}
\end{equation}

陰関数定理に必要な2つの偏導関数を求める。

E9 を m で偏微分して

\begin{equation}
  \frac{\partial F}{\partial m} = a
  \label{eq:E10}
\end{equation}

E9 を $x_1$ で偏微分して

\begin{equation}
  \frac{\partial F}{\partial x_1} = -p_{1} \left(b+a\right)
  \label{eq:E11}
\end{equation}

陰関数定理 $dx_1/dm=-F_m/F_{x_1}$ を研究者が書く。式番号参照で前段の結果を差し込む。

$\frac{dx_1}{dm} = -\frac{E10}{E11}$

\begin{equation}
  \frac{d}{d m} x_{1}=\frac{a}{p_{1} \left(b+a\right)}
  \label{eq:E12}
\end{equation}

一階条件の連立方程式を解いて需要関数を得る。

E6 と E7 と E8 の連立方程式を $x_1$, $x_2$, $\lambda$ について解いて

\begin{equation}
  x_{1} = \frac{a m}{p_{1} b+p_{1} a}
  \label{eq:E13}
\end{equation}

\begin{equation}
  x_{2} = \frac{b m}{p_{2} b+p_{2} a}
  \label{eq:E14}
\end{equation}

\begin{equation}
  \lambda = {{b \left(\frac{a m}{p_{1} b+p_{1} a}\right)^{a} \left(\frac{b m}{p_{2} b+p_{2} a}\right)^{b}+a \left(\frac{a m}{p_{1} b+p_{1} a}\right)^{a} \left(\frac{b m}{p_{2} b+p_{2} a}\right)^{b}}\over{m}}
  \label{eq:E15}
\end{equation}

得られた需要関数を $m$ で直接微分し、陰関数定理から得た感応度と一致することを確かめる。マクロは答えを1つ返すだけで、この比較ができなかった。

E13 を m で偏微分して

\begin{equation}
  \frac{\partial x_1}{\partial m} = \frac{a}{p_{1} b+p_{1} a}
  \label{eq:E16}
\end{equation}

\end{document}